% ==============================================================================
% 00b_fachbegriffe.tex – Verzeichnis der Fachbegriffe (Englisch – Deutsch)
% Wird in 00_vorspann.tex und in Abgabe.tex eingebunden.
% ==============================================================================
\chapter*{Verzeichnis der Fachbegriffe}
\addcontentsline{toc}{chapter}{Verzeichnis der Fachbegriffe}
In dieser Arbeit werden die üblichen englischen Fachbegriffe verwendet. Die
Tabelle nennt die deutsche Übersetzung und eine kurze Erklärung.

{\small\renewcommand{\arraystretch}{1.08}
\begin{longtable}{@{}L{4.1cm}L{3.8cm}>{\raggedright\arraybackslash}p{\dimexpr\textwidth-7.9cm-4\tabcolsep\relax}@{}}
  \toprule
  \textbf{Englisch} & \textbf{Deutsch} & \textbf{Kurz erklärt} \\
  \midrule
  \endhead
  \bottomrule
  \endfoot
  Accuracy & Genauigkeit & systematische Abweichung vom wahren Blickpunkt \\
  Activation (Function) & Aktivierung(sfunktion) & Funktion am Ausgang einer Schicht, z.\,B. Sigmoid \\
  Backbone & Grundnetz & großer vortrainierter Teil des Modells, der die Features berechnet \\
  Baseline & Vergleichswert & einfaches Vergleichsverfahren, z.\,B. immer derselbe Punkt \\
  Batch Size & Batchgröße & Anzahl der Bilder pro Trainingsschritt \\
  Chicken-and-Egg Problem & Henne-Ei-Problem & zwei Dinge hängen gegenseitig voneinander ab \\
  Confidence Interval (CI) & Konfidenzintervall & Bereich, in dem der wahre Mittelwert mit hoher Wahrscheinlichkeit liegt \\
  Crowdsourcing & -- & viele Freiwillige nehmen Daten selbst auf \\
  Data Leakage & Datenleck & Information aus den Testdaten gelangt ins Training \\
  Data Loss & Datenverlust & Anteil der Messungen ohne gültigen Wert \\
  Domain Adaptation & Domänenadaption & Anpassung an eine neue Umgebung mit Daten von dort \\
  Domain Generalization & Domänengeneralisierung & gute Leistung in neuen Umgebungen ohne Daten von dort \\
  Domain Shift & Domänenwechsel & Testdaten unterscheiden sich von den Trainingsdaten \\
  Early Stopping & vorzeitiges Beenden & Training endet, wenn sich die Validierung nicht mehr verbessert \\
  Epoch & Epoche & ein Durchlauf durch alle Trainingsdaten \\
  Feature(s) & Merkmal(e) & Zahlen, die ein Netz aus einem Bild berechnet \\
  Feature Extractor & Merkmalsextraktor & Netz, das die Features berechnet \\
  Feature Space & Merkmalsraum & Raum, in dem die Features liegen \\
  Fine-Tuning & Feinabstimmung & ein vortrainiertes Modell mit neuen Daten weitertrainieren \\
  Gaze Regression & Blickregression & Blick als kontinuierlicher Wert vorhersagen \\
  Head-Training & -- & nur der Regression Head wird trainiert \\
  Image / Text Encoder & Bild- / Textencoder & Netz, das ein Bild bzw. einen Text in Features umwandelt \\
  Label Noise & Label-Rauschen & Labels stimmen nicht genau mit dem echten Blick überein \\
  Learning Rate & Lernrate & Schrittweite beim Training \\
  Loss / Loss Function & Verlust / Verlustfunktion & Maß für den Fehler, das beim Training minimiert wird \\
  Overfitting & Überanpassung & Modell lernt die Trainingsdaten auswendig und wird auf neuen Daten schlechter \\
  Precision & Präzision & Streuung wiederholter Messungen \\
  Pretraining & Vortraining & Training auf sehr vielen allgemeinen Daten vor der eigentlichen Aufgabe \\
  Regression Head & Regressionskopf & kleines Netz, das aus den Features die Koordinaten berechnet \\
  Remote Camera & entfernte Kamera & Kamera steht vor der Person, nicht am Kopf \\
  Ridge Regression & Ridge-Regression & lineare Regression mit Strafe für große Gewichte \\
  Saccade & Sakkade & sehr schnelle Augenbewegung zu einem neuen Ziel \\
  Self-Supervised Learning & selbstüberwachtes Lernen & Lernziel entsteht aus den Daten selbst, ohne Labels \\
  Split & Aufteilung & Aufteilung der Daten in Training, Validierung und Test \\
  Supervised Learning & überwachtes Lernen & Lernen mit Labels, die Menschen vergeben haben \\
  Test-Time Adaptation (TTA) & Anpassung zur Testzeit & Anpassung erst während der Nutzung \\
  Transition Frames / Period & Übergangsbilder / Übergangszeit & Bilder bzw. Zeit direkt nach einem Zielsprung \\
  Unconstrained Conditions & uneingeschränkte Bedingungen & freie Kopfhaltung, wechselndes Licht, beliebiger Abstand \\
  Unsupervised Domain Adaptation & unüberwachte Domänenadaption & Domain Adaptation ohne Labels der Zielumgebung \\
  Weakly Supervised & schwach überwacht & Lernen mit wenigen oder groben Labels \\
  Weight Decay & Gewichtsabnahme & Strafe für große Gewichte gegen Overfitting \\
\end{longtable}}
